% ============================================
% Title Options
% ============================================

\title{\revsecond{Interspecies interaction strength affects community-level selection in microbial coalescence}}

% Title revised at the editor's request, second revision round.
% Previous title: Interspecies Interactions Drive Community-Level Selection in Microbial Coalescence

% Alternative title:
% Interaction strength determines when microbial communities behave as units of selection

% ============================================
% Authors and Affiliations
% ============================================
\author[1]{Jinyeop Song}
\author[1,2]{Jiliang Hu}
\author[1]{Jeff Gore}

\affil[1]{\revsecond{Physics of Living Systems, }Department of Physics, Massachusetts Institute of Technology, Cambridge, MA, USA}
\affil[2]{\revsecond{Department of Mechanical Engineering, Massachusetts Institute of Technology, Cambridge, MA, USA}}

\date{}

\maketitle

% ============================================
% Abstract
% ============================================
\begin{abstract}
\revsecond{Whether communities behave as cohesive units or as loose collections of independent species is a fundamental question in ecology.} Here, we study this question in the context of community coalescence, the mixing of previously isolated communities, using \revsecond{synthetic} bacterial microcosm experiments combined with \revsecond{Lotka--Volterra} modeling. Our results demonstrate that effective interaction intensity and environmental feedbacks determine whether communities or species are the units of selection during coalescence. When effective interactions are moderate to strong, one parental community consistently outcompetes the other, indicating community-level selection. In contrast, under weak interactions, species fates are uncorrelated and the two communities contribute equally to the coalesced outcome, indicating the absence of community-level selection. \rev{A similar nutrient-dependent shift toward \revsecond{dominance of one parental community} also appears in taxonomically richer communities derived from natural environmental samples.} Furthermore, we identify two distinct regimes underlying community-level selection in experiments with different media conditions: an emergent regime in which collective dynamics shape outcomes that cannot be predicted from species traits alone, and a top-down regime where dominant species determine the winning community. Together, these results reconcile conflicting observations on community-level selection during community coalescence by demonstrating that communities behave as cohesive units only when effective interactions and feedbacks are sufficiently strong.
\end{abstract}

\newpage
